\section{Compact Group Scan and Finite-Resolution Projection Readout}

\subsection{Scanning Dynamics on Compact Groups}

Let $G$ be a compact group; the most basic case takes $G=\TT^d=(\RR/\ZZ)^d$. Given a scanning step vector $\alpha\in\TT^d$, define the discrete scanning orbit
\[
x_n = x_0 + n\alpha \pmod{1},\qquad n\in\NN.
\]
When the coordinates of $\alpha$ are linearly independent over $\QQ$, the orbit is dense and uniquely ergodic on $\TT^d$, so time averages equal space averages for continuous functions. This property forms the mathematical foundation of the ``scan---average'' chain.

\subsection{Finite-Resolution Readout and Window Kernels}

A finite observer cannot access the full information of continuous phase; we must introduce a finite-resolution instrument. Let $\Pi$ be the spectral measure of phase (or in the classical case a Borel measure), and take a family of window kernels $w_k^{(\varepsilon)}$ (where $\varepsilon$ denotes resolution). Define effect operators
\[
E_k^{(\varepsilon)}=\int w_k^{(\varepsilon)}(x)\,\dd\Pi(x),
\]
forming a POVM. For binary readout, we can take $k\in\{0,1\}$ and let window $W\subset \TT$, defining via indicator function or smooth approximation
\[
s_n = \ind_W(x_n)\in\{0,1\}.
\]
Using Lebesgue measure as reference, under unique ergodicity we have $\lim_{N\to\infty}\frac1N\sum_{n=0}^{N-1}s_n=\mu(W)$. Finite $N$ deviation is controlled by discrepancy.

\subsection{Discrepancy Certificate and Auditable Error Bounds}

Define one-dimensional star discrepancy
\[
D_N^*=\sup_{a\in[0,1]}\left|\frac1N\#\{0\le n<N:x_n<a\}-a\right|.
\]
For bounded variation functions $f$, Koksma--Hlawka-type inequalities give deterministic error bounds
\[
\left|\frac1N\sum_{n=0}^{N-1}f(x_n)-\int_0^1 f(x)\,\dd x\right|\le \Var(f)\,D_N^*.
\]
Therefore, if we write readouts or statistics as time averages of phase functions, we can obtain ``auditable'' finite-sample error certificates by computing/estimating $D_N^*$, without resorting to probabilistic limit theorems.

\subsection{Weyl Pairs and Intrinsic Complementarity}

In the operator formulation, scanning is realized by a unitary operator $U_{\mathrm{scan}}$ (e.g., $(U_{\mathrm{scan}}\psi)(x)=\psi(x+\alpha)$), and pointer phase by a multiplication operator $V$ (e.g., $(V\psi)(x)=e^{2\pi i x}\psi(x)$). The two satisfy the Weyl relation
\[
U_{\mathrm{scan}}V=e^{2\pi i\alpha}VU_{\mathrm{scan}}.
\]
When $\alpha\notin\QQ$, they have no nonzero common eigenvector, thus providing the structural reason why scanning time and pointer phase cannot be simultaneously sharp. This complementarity provides the foundation for subsequent unified treatment of ``readout probability---irreversibility---resource cost.''
